\documentclass[11pt]{article}
\usepackage[localtheorems]{raeez-math-template}

%%% Packages

%%% Page geometry

%%% Theorem environments
\theoremstyle{plain}
\newtheorem{theorem}{Theorem}[section]
\newtheorem{proposition}[theorem]{Proposition}
\newtheorem{corollary}[theorem]{Corollary}
\newtheorem{lemma}[theorem]{Lemma}
\newtheorem{conjecture}[theorem]{Conjecture}

\theoremstyle{definition}
\newtheorem{definition}[theorem]{Definition}
\newtheorem{construction}[theorem]{Construction}
\newtheorem{example}[theorem]{Example}

\theoremstyle{remark}
\newtheorem{remark}[theorem]{Remark}

%%% Macros
\newcommand{\fg}{\mathfrak{g}}
\newcommand{\fh}{\mathfrak{h}}
\newcommand{\fgl}{\mathfrak{gl}}
\newcommand{\fsl}{\mathfrak{sl}}
\newcommand{\fL}{\mathfrak{L}}
\newcommand{\Ran}{\mathrm{Ran}}
\newcommand{\Mbar}{\overline{\mathcal{M}}}
\newcommand{\FM}{\mathrm{FM}}
\newcommand{\MC}{\mathrm{MC}}
\newcommand{\cA}{\mathcal{A}}
\newcommand{\cC}{\mathcal{C}}
\newcommand{\cH}{\mathcal{H}}
\newcommand{\cM}{\mathcal{M}}
\newcommand{\cO}{\mathcal{O}}
\newcommand{\cW}{\mathcal{W}}
\newcommand{\cZ}{\mathcal{Z}}
\newcommand{\bC}{\mathbb{C}}
\newcommand{\bP}{\mathbb{P}}
\newcommand{\bR}{\mathbb{R}}
\newcommand{\bS}{\mathbb{S}}
\newcommand{\bZ}{\mathbb{Z}}
\newcommand{\Ainf}{A_\infty}
\newcommand{\Eone}{E_1}
\newcommand{\Etwo}{E_2}
\newcommand{\Ethree}{E_3}
\newcommand{\En}{E_n}
\newcommand{\Sd}{\bS^d}
\newcommand{\Res}{\operatorname{Res}}
\newcommand{\Tr}{\operatorname{Tr}}
\newcommand{\HH}{\mathrm{HH}}
\newcommand{\HC}{\mathrm{HC}}
\newcommand{\CC}{\mathrm{CC}}
\newcommand{\Zder}{Z^{\mathrm{der}}_{\mathrm{ch}}}
\newcommand{\ChirHoch}{\mathrm{ChirHoch}}
\newcommand{\CoHA}{\mathrm{CoHA}}
\newcommand{\Fuk}{\mathrm{Fuk}}
\newcommand{\Rep}{\mathrm{Rep}}
\newcommand{\Ext}{\mathrm{Ext}}
\newcommand{\Fact}{\mathrm{Fact}}
\newcommand{\CY}{\mathrm{CY}}
\newcommand{\DT}{\mathrm{DT}}
\newcommand{\GW}{\mathrm{GW}}
\newcommand{\Hom}{\mathrm{Hom}}
\newcommand{\Conf}{\mathrm{Conf}}
\newcommand{\MF}{\mathrm{MF}}
\newcommand{\Coh}{\mathrm{Coh}}
\newcommand{\SCchtop}{\mathrm{SC}^{\mathrm{ch},\mathrm{top}}}
\newcommand{\hCS}{\mathrm{hCS}}
\newcommand{\BRST}{\mathrm{BRST}}
\newcommand{\Uq}{U_q}
\DeclareMathOperator{\Aut}{Aut}
\DeclareMathOperator{\Spec}{Spec}

%%% Title
\title{Calabi--Yau quantum groups\\
and $6$d holomorphic Chern--Simons}
\author{Raeez Lorgat}
\date{}

\begin{document}
\maketitle

%%% ================================================================
%%% ABSTRACT
%%% ================================================================

\begin{abstract}
Volume~I of the modular Koszul duality programme introduces
$\Eone$-chiral quantum groups as abstract algebraic-geometric
objects: ordered bar complexes with spectral $R$-matrices, Drinfeld
coproducts, and Yang--Baxter equations on configuration spaces of
curves.  This paper constructs them concretely from Calabi--Yau
geometry.

The source is the holomorphic Chern--Simons hierarchy.  At the top
($d = 6$): holomorphic Chern--Simons theory on a CY$_3$-fold~$Y$.
A curve $C \subset Y$ of codimension~$2$ carries a defect algebra:
the normal bundle $N_{C/Y}$ supplies two spectral parameters, and
the defect algebra is a chiral algebra on~$C$ with quantum group
structure from the normal directions.  For $Y = \bC^3$ and
$C = \bC$, the local defect gives the abelian low-spin witness
$J(z)J(w)\sim\Psi_{\mathrm{lev}}/(z-w)^2$, with
$\Psi_{\mathrm{lev}}=-\sigma_2(h_1,h_2,h_3)$ and Sugawara
central charge~$1$.  The full $\cW_{1+\infty}$ object requires the
completed Drinfeld double or Fock/evaluation data. The hierarchy
$6 \to 5 \to 4 \to 3$ is codimension-$2$ restriction followed by
curve specialization; the ambient hCS observables are holomorphic
$E_3$ in complex dimension~$3$ and holomorphic $E_2$ in complex
dimension~$2$, while the curve output is $\Eone$.

The critical cohomological Hall algebra $\CoHA(Q, W)$ of a quiver
with CY$_3$ potential is the positive $\Eone$-primitive: its product
is ordered Hall multiplication. The full Hopf and vertex-coproduct
data require localization, a Hopf pairing, the opposite half, the
Cartan sector, and completion. For the Jordan quiver:
	$\CoHA(\bC^3) \simeq Y^+(\widehat{\fgl}_1)$. The full
	$\cW_{1+\infty}$ vertex realisation requires the corresponding
	double/evaluation construction.
The Drinfeld center recovers $\Etwo$-braiding, closing the $\En$
operadic circle.

Quantum toroidal algebras
$U_{q,t}(\widehat{\widehat{\fgl}}_1)$ carry the CY condition as
the parameter constraint $q_1 q_2 q_3 = 1$.  The structure
function
$G(x; q_1, q_2, q_3) = \prod_i(1 - q_i x)/\prod_i(1 - q_i^{-1}x)$
satisfies $G(x)G(1/x)=(q_1q_2q_3)^2$; the identity
$G(x)G(1/x)=1$ characterizes the connected CY branch
$q_1q_2q_3=1$ after the exponential/formal branch is fixed.

Wall-crossing in the Bridgeland stability manifold is
Maurer--Cartan gauge equivalence in the convolution algebra of
Volume~I: the Kontsevich--Soibelman wall-crossing formula is the
MC gauge transformation $\Theta_\sigma \sim \Theta_{\sigma'}$.
The shadow obstruction tower at $\kappa_{\mathrm{ch}} = \Psi$
recovers the genus-$1$ constant-map contribution for $\bC^3$; the
all-genus virtual equality is conditional on the uniform-weight
and $\Phi_3$ hypotheses, and the motivic equality remains
conjectural in general.
\end{abstract}

\tableofcontents

%%% ================================================================
%%% 1. INTRODUCTION
%%% ================================================================

\section{Introduction}\label{sec:intro}

\subsection{CY quantum groups as concrete $\Eone$-chiral quantum
groups}

Volume~I constructs $\Eone$-chiral quantum groups as abstract
algebraic-geometric objects on curves.  An $\Eone$-chiral quantum
group is a tuple $(A, \mu, \Delta_z, \varepsilon, \eta, S)$
satisfying five axioms (H1--H5): an $\Eone$-chiral algebra
structure, a spectral coproduct, bialgebra compatibility, spectral
coassociativity
\[
 (\Delta_w \otimes \mathrm{id}) \circ \Delta_{z+w}
 = (\mathrm{id} \otimes \Delta_z) \circ \Delta_w,
\]
and the Hopf axiom at $z = 0$.  The ordered bar
$B^{\mathrm{ord}}(A) = T^c(s^{-1}\bar{A})$ preserves the
$R$-matrix; the symmetric bar $B^\Sigma(A)$ kills the Hopf
structure via the lossy averaging map: in abelian and scalar
families $\mathrm{av}(r(z)) = \kappa_{\mathrm{ch}}$, while for
non-abelian affine Kac--Moody
$\mathrm{av}(r(z))+\dim(\fg)/2 = \kappa_{\mathrm{ch}}$.

Two geometric inputs produce concrete $\Eone$-chiral quantum groups:
the $6$d holomorphic Chern--Simons hierarchy, which gives defect
chiral algebras on codimension-$2$ curves, and critical CoHA, which
gives the ordered $\Eone$-primitive graded by dimension vectors.

The $\Eone$ stabilization theorem of the companion paper
proves that for $d \geq 3$ the constructed CY chiral output on the
reference curve is natively $\Eone$ on the witnessed loci.  The full
Schouten--Nijenhuis bracket is nonzero; the ordered CoHA direction is
the quantum-group source.  The braided $\Etwo$-structure is recovered
through the Drinfeld center of the $\Eone$-monoidal representation
category. For $d\ge3$,
\[
\pi_1(\Conf_2(\bR^d))=\pi_1(S^{d-1})=0,
\qquad
\pi_1(\Conf_2(\bR^d)/\mathfrak S_2)=\bZ/2,
\]
so the unordered quotient supplies only symmetric braiding.

\subsection{The holomorphic Chern--Simons hierarchy}

The theories of this paper form a hierarchy in real dimension:
\begin{center}
\renewcommand{\arraystretch}{1.3}
\small
\begin{tabular}{llll}
 \toprule
 $d$ & Theory & Curve output & $\En$ level \\
 \midrule
 $6$ & hCS on CY$_3$ $Y$ & defect on $C \subset Y$ &
 $\Eone$ \\
 $5$ & partially topological CS on $\bR \times S$ &
 $\bR \times C$ & $\Eone$ \\
 $4$ & hCS on $\bC^2$ (CWY) & line operator & $\Eone$ \\
 $3$ & 3d HT (Volume~II) & boundary algebra & $\Eone$ \\
 \bottomrule
\end{tabular}
\end{center}
The hierarchy $6 \to 5 \to 4 \to 3$ is codimension-$2$ restriction
at each step.  Each restriction produces a chiral algebra from a
gauge theory.  The resulting chiral algebras are $\Eone$-chiral
quantum groups in the sense of Volume~I, with $R$-matrices from the
normal directions and coproducts from the spectral parameters.

\subsection{The $\En$ operadic circle}

The $\En$ hierarchy organizes the entire programme:
\[
 \Ethree(\text{bulk}) \;\to\;
 \Etwo(\text{boundary chiral}) \;\to\;
 \Eone(\text{bar/QG}) \;\to\;
 \Etwo(\text{Drinfeld center}) \;\to\;
 \Ethree(\text{derived center}).
\]
Each arrow has precise mathematical content:
$\Ethree \to \Etwo$ is restriction to a codimension-$2$ defect;
$\Etwo \to \Eone$ is the ordered bar complex;
$\Eone \to \Etwo$ is the Drinfeld center;
$\Etwo \to \Ethree$ is the higher Deligne conjecture.  The circle
closes for 3d HT theories with conformal vector at non-critical
level (the topologization theorem, proved for affine KM); without
conformal vector, the theory is stuck at $\SCchtop$.

The CY quantum groups of this paper enter the circle at the
$\Eone$ vertex: the CoHA is the $\Eone$-primitive, and the
Drinfeld center carries it to $\Etwo$.  The tetrahedron equation
(3d analogue of YBE) fails for the Yang $R$-matrix at
$O(\kappa_{\mathrm{ch}}^2)$: the $\Ethree$ structure is
genuinely nontrivial beyond $\Etwo$.

\subsection{Construction layers}

The construction has five layers.  Codimension-$2$ hCS restriction
produces the defect chiral algebra on a curve.  The local $\bC^3$
defect identifies the Heisenberg low-spin channel and isolates the
completed double/Fock data needed for $\cW_{1+\infty}$.  Critical
CoHA supplies the ordered positive half.  The toroidal formal disk
imposes $q_1q_2q_3=1$.  The Drinfeld center and the shadow tower
connect the $\Eone$ primitive to $\Etwo$ braiding and the genus-one
DT/GW channel.


%%% ================================================================
%%% 2. THE hCS HIERARCHY
%%% ================================================================

\section{The holomorphic Chern--Simons hierarchy}
\label{sec:hcs-hierarchy}

\subsection{$6$d holomorphic Chern--Simons}

Let $Y$ be a CY$_3$-fold with holomorphic volume form
$\Omega_Y \in H^0(Y, \omega_Y)$.  The holomorphic Chern--Simons
functional on $Y$ (Witten, Costello) is
\[
 S_{\hCS}(\cA)
 = \int_Y \Omega_Y \wedge
 \mathrm{tr}\Bigl(\cA \wedge \bar\partial \cA
 + \tfrac{2}{3}\cA \wedge \cA \wedge \cA\Bigr)
\]
where $\cA \in \Omega^{0,1}(Y, \fg)$ is a $(0,1)$-form valued in
a Lie algebra~$\fg$.  The equations of motion
$\bar\partial \cA + \cA \wedge \cA = 0$ are the integrability
condition for the holomorphic structure on the trivial $G$-bundle.

The theory is six-dimensional, holomorphic in all three complex
directions.  The BRST operator is $Q_{\BRST} = \bar\partial +
[\cA, -]$, and the observables are the Dolbeault cohomology
$H^{0,\bullet}_{\bar\partial_\cA}(Y, \fg)$.

\begin{remark}[Relation to physical dimensions]
\label{rem:v3-qg-hcs-physical}
The $6$d hCS theory is the Costello--Li/BCOV open B-model
holomorphic field-theory avatar of the $E_3$ story: the kinetic term
is first-order, the path integral is one-loop exact, and the partition
function is a holomorphic invariant.  The $6$d $(2,0)$ Schur/BLLPR
route is an independent construction.  The physical content sits in
the defect algebras.
\end{remark}

\subsection{Dimensional reduction: $6 \to 5 \to 4 \to 3$}

The hierarchy proceeds by codimension-$2$ restriction at each step.

\emph{$6 \to 5$}: for a CY$_3$-fold of the form
$Y = \bR \times S$ (where $S$ is a CY$_2$ surface), the $6$d hCS
theory reduces to $5$d partially topological Chern--Simons on
$\bR \times S$.  The topological direction $\bR$ carries the
$\Eone$ ordering; the holomorphic surface $S$ carries the chiral
data.  Restriction to $\bR \times C$ (where $C \subset S$ is a
curve) produces the ordered bar complex with one spectral parameter.

\emph{$5 \to 4$}: for $S = \bC^2$, the $5$d theory reduces to
$4$d holomorphic Chern--Simons on $\bC^2$: the
Costello--Witten--Yamazaki theory.  Line operators on $\bC \subset
\bC^2$ produce Yangians from holomorphic gauge fields.  The
$R$-matrix comes from the normal direction to the line.

\emph{$4 \to 3$}: restricting to a further codimension-$2$ locus
produces the $3$d holomorphic-topological theories of Volume~II.
The boundary algebra is a chiral algebra on the codimension-$2$
curve, and the $\SCchtop$ pair
$(C^\bullet_{\mathrm{ch}}(A,A), A)$ is the boundary--bulk system.

\begin{proposition}[Dimensional consistency]
\label{prop:v3-qg-dim-consistency}
The ambient hCS observables retain their holomorphic factorization
level before curve specialization: $6$d hCS on a CY$_3$ has ambient
holomorphic $E_3$ observables, and $5$d hCS on $\bR\times\bC^2$
has holomorphic $E_2$ observables. Projection to a defect curve is
the stage that produces the ordered $\Eone$ chiral algebra with
spectral parameters from the normal directions. The $\Etwo$-braiding
is recovered through the Drinfeld center.
\end{proposition}

\subsection{The $\Omega$-background and equivariant parameters}

For the toric CY$_3$-fold $Y = \bC^3$, the three coordinate
directions carry equivariant parameters
$h_1, h_2, h_3 \in \bC$ subject to the CY condition
$h_1 + h_2 + h_3 = 0$.  In the multiplicative parametrization
$q_i = e^{h_i}$, the CY condition reads $q_1 q_2 q_3 = 1$.  The
identification $q_1 = q$, $q_2 = t^{-1}$,
$q_3 = t/q$ recovers the standard $(q,t)$-parametrization of the
quantum toroidal algebra.

The $\Omega$-background parameters enter the defect algebra through
the equivariant Chern--Simons functional: the gauge field on $Y$
decomposes along the normal bundle to a defect curve
$C \subset Y$, and the equivariant parameters weight the normal
modes.  For $C = \bC_z$ (the first coordinate axis) inside
$Y = \bC^3$, the normal bundle is
$N_{C/Y} = \cO_C^{\oplus 2}$ with fiber coordinates
$(w_1, w_2)$ carrying equivariant weights $(h_2, h_3)$.  The
resulting defect algebra on $C$ has OPE poles involving both normal
parameters.


%%% ================================================================
%%% 3. CODIMENSION-2 DEFECT ALGEBRAS
%%% ================================================================

\section{Codimension-$2$ defect algebras}
\label{sec:codim2}

\subsection{The defect construction}

A curve $C \subset Y$ in a CY$_3$-fold carries a defect algebra:
the $6$d hCS theory restricted to a tubular neighborhood of $C$
produces a $2$d chiral algebra on $C$ with additional structure from
the normal bundle.

\begin{construction}[Codimension-$2$ defect algebra]
\label{con:v3-qg-codim2}
Let $Y$ be a CY$_3$-fold, $C \subset Y$ a smooth curve, and
$N_{C/Y}$ the normal bundle of rank~$2$.  The defect algebra
$A_C^{\mathrm{def}}$ on $C$ is obtained by:
\begin{enumerate}[label=\textbf{\arabic*.}]
 \item \textbf{Restriction.}  Restrict the $6$d hCS gauge field
 $\cA \in \Omega^{0,1}(Y, \fg)$ to the formal neighborhood
 $\hat{N}_{C/Y}$ of $C$ in $Y$.  The Taylor expansion along the
 normal fibers gives a collection of fields on $C$ indexed by
 normal-direction monomials.
 \item \textbf{BRST reduction.}  The $\BRST$ operator
 $Q = \bar\partial + [\cA, -]$ restricts to the formal
 neighborhood.  The $Q$-cohomology in the fiber directions
 produces the generating fields of the defect algebra on~$C$.
 \item \textbf{OPE from normal modes.}  The OPE on $C$ is
 inherited from the $6$d hCS action.  The normal bundle
 $N_{C/Y}$ supplies two spectral parameters (from the two fiber
 coordinates), which enter the OPE as meromorphic functions of
 the insertion points on~$C$.
\end{enumerate}
\end{construction}

\begin{remark}[Two spectral parameters]
\label{rem:v3-qg-two-spectral}
The rank-$2$ normal bundle $N_{C/Y}$ gives two spectral parameters.
This is the geometric origin of the two-parameter family of quantum
toroidal algebras: the two normal directions to a curve inside a
CY$_3$-fold produce the $(q,t)$-deformation.  For a CY$_2$-fold
(surface), the normal bundle has rank~$1$, producing one spectral
parameter: the standard spectral parameter of Yangians and
KZ equations.
\end{remark}

\subsection{$\bC^3$: the low-spin defect witness}

\begin{proposition}[Low-spin defect algebra of $\bC^3$]
\label{prop:v3-qg-c3-defect}
For $Y = \bC^3$ with coordinates $(z, w_1, w_2)$ and CY form
$\Omega = dz \wedge dw_1 \wedge dw_2$, the codimension-$2$ defect
algebra on $C = \bC_z$ has a Heisenberg current
\[
J(z)J(w)\sim \frac{\Psi_{\mathrm{lev}}}{(z-w)^2},
\qquad
\Psi_{\mathrm{lev}}=-\sigma_2(h_1,h_2,h_3)
=\frac{h_1^2+h_2^2+h_3^2}{2}
\]
on the Calabi--Yau branch $h_1+h_2+h_3=0$. Its Sugawara stress
tensor is
\[
T=\frac{1}{2\Psi_{\mathrm{lev}}}{:}J^2{:},\qquad c=1.
\]
The all-spin $\cW_{1+\infty}$ vertex algebra requires the completed
Drinfeld double/Fock evaluation data in addition to this local
abelian witness.
\end{proposition}

\begin{proof}
The gauge field $\cA \in \Omega^{0,1}(\bC^3, \fgl_N)$ expands along
the normal fiber as
$\cA(z, w_1, w_2) = \sum_{a,b \geq 0}
J_{a,b}(z)\, w_1^a w_2^b\, d\bar{z}$,
where $J_{a,b}$ are fields on $C = \bC_z$.  The lowest BRST class is
the $\fgl_1$ current $J(z)$. The OPE is computed from the quadratic
part of the $6$d hCS action restricted to the formal neighborhood,
and the normal equivariant weights enter through
$\Psi_{\mathrm{lev}}=-\sigma_2(h_1,h_2,h_3)$. The CY condition
$h_1 + h_2 + h_3 = 0$ gives
$\sigma_2 = -(h_1^2 + h_2^2 + h_3^2)/2$, so
$\Psi_{\mathrm{lev}}=(h_1^2+h_2^2+h_3^2)/2$. The Sugawara formula
for a nonzero-level Heisenberg current gives $c=1$.
\end{proof}

\begin{remark}[Level-prefixed $r$-matrix]
\label{rem:v3-qg-level-r-matrix}
The classical $r$-matrix of the $\widehat\fgl_1$ current subalgebra
derived from the 6d reduction carries level $\Psi = -\sigma_2$, not
a normalised level $k = 1$:
\[
 r^{\mathrm{Heis}}(z) = \frac{\Psi}{z}
 = \frac{-\sigma_2(h_1,h_2,h_3)}{z}.
\]
At $\Psi = 0$ the $r$-matrix vanishes and the Sugawara construction
is singular; this is the degenerate locus sitting outside the
non-critical Koszul locus.  The averaging map gives
$\mathrm{av}(r(z)) = \Psi = \kappa_{\mathrm{ch}}(H_\Psi)$.  For the
full $\cW_{1+\infty}$ algebra, higher-spin $r$-matrix components
enter through the universal Miura cross-term: the leading term
$J \otimes W_{s-1}$ has coefficient $(\Psi - 1)/\Psi$ for every
$s \geq 2$. The central charge $c = 1$ is the rank of the underlying
free boson and is independent of $\Psi$; Prochazka--Rapcak
parametrise $\cW_{1+\infty}$ by
$(\lambda_1,\lambda_2,\lambda_3)$ with
$\sum_i 1/\lambda_i=0$. The normalisation $\Psi=-\sigma_2$ is the one
for which the equivariant hCS quadratic form gives
$J(z)J(w)\sim\Psi/(z-w)^2$.
\end{remark}

\subsection{Defect algebras for general CY$_3$}

For a general CY$_3$-fold $Y$ with a curve $C \subset Y$, the
defect algebra depends on the normal bundle $N_{C/Y}$.  The normal
bundle is a rank-$2$ vector bundle on $C$; its splitting type and
degree control the defect algebra.

\begin{example}[Conifold]
\label{ex:v3-qg-conifold-defect}
For the resolved conifold $Y = \mathrm{Tot}(\cO(-1) \oplus
\cO(-1) \to \bP^1)$ and $C = \bP^1$ (the zero section):
$N_{C/Y} = \cO(-1)^{\oplus 2}$.  The defect algebra on $C$ has
$\kappa_{\mathrm{ch}} = 1$ (from the DT computation) and shadow
class~G (shadow depth~$2$).  The negative degrees of the normal
bundle constrain the spin content: only finitely many spins
contribute, and the resulting algebra is a truncation of
$\cW_{1+\infty}$.
\end{example}

\begin{example}[Local $\bP^2$]
\label{ex:v3-qg-local-p2-defect}
For $Y = \mathrm{Tot}(K_{\bP^2} \to \bP^2)$ and a line
$C = \bP^1 \subset \bP^2$:
$N_{C/Y} = \cO(1) \oplus \cO(-3)$.  The defect algebra has
$\kappa_{\mathrm{ch}} = 3/2$ and shadow class~M (infinite shadow
depth).  The mixed splitting type of the normal bundle introduces
asymmetry between the two spectral parameters.
\end{example}


%%% ================================================================
%%% 4. CoHA AS E1 PRIMITIVE
%%% ================================================================

\section{The CoHA as $\Eone$-primitive}
\label{sec:coha-e1}

\subsection{Critical cohomological Hall algebras}

For a quiver with potential $(Q, W)$ arising from a toric
CY$_3$-fold, the critical CoHA is the $\bZ$-graded associative
algebra
\[
 \CoHA(Q, W)
 = \bigoplus_{\mathbf{d} \in \bZ_{\geq 0}^{Q_0}}
 H^{\mathrm{BM}}_*(\mathrm{Crit}(W_\mathbf{d}),
 \phi_{W_\mathbf{d}})
\]
where $\phi_{W_\mathbf{d}}$ is the vanishing cycle sheaf and the
multiplication comes from Hall-algebra extension of representations.
The grading by dimension vectors $\mathbf{d}$ is the CY analogue of
the weight-space decomposition of a quantum group.

\begin{theorem}[Schiffmann--Vasserot; Rapcak--Soibelman--Yang--Zhao]
\label{thm:v3-qg-coha-yangian}
For a toric CY$_3$ $X$ without compact $4$-cycles:
\[
 \CoHA(Q_X, W_X) \;\simeq\; Y^+(\widehat{\fg}_{Q_X})
\]
where $Y^+(\widehat{\fg}_{Q_X})$ is the positive half of the
affine super Yangian attached to the toric quiver.  For $\bC^3$
(Jordan quiver, cubic potential):
$\CoHA(\bC^3) \simeq Y^+(\widehat{\fgl}_1)$. The
$\cW_{1+\infty}$ vertex realisation appears only after passing to
the Drinfeld double, or to the corresponding evaluation shadow.
\end{theorem}

\subsection{The CoHA carries the positive $\Eone$ half}

The bare critical CoHA is the positive half of the CY quantum group.
It supplies the ordered Hall product. Spectral coproducts, Hopf
pairings, antipodes, Cartan currents, and negative roots belong only
after localization, completion, and passage to the Drinfeld double.

\begin{proposition}[CoHA positive-half data]
\label{prop:v3-qg-coha-hopf}
The critical CoHA $\CoHA(Q, W)$ of a toric CY$_3$ carries the
positive $\Eone$ data:
\begin{enumerate}[label=\textup{(H\arabic*)}]
 \item \emph{Ordered Hall product.}  The Hall multiplication
 (extension of short exact sequences $0 \to V' \to V \to V'' \to 0$
 with sub before quotient) is an ordered ($\Eone$) product.  The
 ordering is physical: the dimension-vector filtration imposes a
 preferred direction.
 \item \emph{Positive half.}  For the Jordan quiver, this ordered
 product identifies $\CoHA(\bC^3)$ with
 $Y^+(\widehat{\fgl}_1)$.
 \item \emph{Double input.}  A spectral coproduct or Hopf antipode
 requires the opposite half, Cartan sector, Hopf pairing, and
 completed Drinfeld double. These structures are extra data beyond
 the bare positive CoHA.
\end{enumerate}
\end{proposition}

\begin{remark}[Hopf data require the double]
\label{rem:v3-qg-antipode}
The transfer-matrix antipode $S(T(u))=T(u)^{-1}$ is a statement in
the completed Yangian or toroidal double. The Hopf identity requires
the negative half and the Cartan current in addition to the positive
CoHA.
\end{remark}

\subsection{The $R$-matrix and the Drinfeld coproduct}

The $R$-matrix of the CoHA arises from the Ext-quiver: the
extension groups $\Ext^1(V', V'')$ between representations
determine the poles of the $R$-matrix, and the CY condition
(self-duality of the Ext-quiver) constrains the $R$-matrix
to satisfy the Yang--Baxter equation.

\begin{proposition}[Universal Miura coefficient]
\label{prop:v3-qg-miura-universal}
The Drinfeld coproduct $\Delta_z$ for $\cW_{1+\infty}$ has the
universal Miura coefficient $(\Psi - 1)/\Psi$ on the leading
cross-term $J \otimes W_{s-1} + W_{s-1} \otimes J$ at every spin
$s \geq 2$.  At spin~$2$:
\[
 \Delta_z(T)
 = T \otimes 1 + 1 \otimes T
 + \frac{\Psi - 1}{\Psi}\,J \otimes J
 + O(z^{-1}).
\]
The coefficient $(\Psi - 1)/\Psi$ is structural: it arises from
the Prochazka--Rapcak Miura factorization $T(u) =
\prod_i (u + \Lambda_i)$ in three steps: (1)~single-$J$ sector
gives $1/\Psi$; (2)~Drinfeld shift adds $+1$; (3)~lower sectors do
not contribute.  The coefficient persists at all spins $s \geq 2$ with
direct symbolic expansion through spins $2$--$6$ giving the same
coefficient.
\end{proposition}

\begin{remark}[Coincidental agreement at $\Psi = 2$]
\label{rem:v3-qg-psi2-coincidence}
At $\Psi = 2$: $(\Psi - 1)/\Psi = 1/2 = 1/\Psi$.  The two
expressions $(\Psi - 1)/\Psi$ and $1/\Psi$ agree at this single
point, which is the self-dual level.  The distinction is visible at
$\Psi = 3$: $(\Psi - 1)/\Psi = 2/3 \neq 1/3 = 1/\Psi$.  Verify
at three or more parameter values to distinguish the two.
\end{remark}


%%% ================================================================
%%% 5. QUANTUM TOROIDAL ALGEBRAS
%%% ================================================================

\section{Quantum toroidal algebras and the CY condition}
\label{sec:toroidal}

The quantum toroidal algebra
$U_{q,t}(\widehat{\widehat{\fgl}}_1)$ is the two-parameter
deformation of the double current algebra $\fgl_1 \otimes
\bC[u^{\pm 1}, v^{\pm 1}]$, carrying the full
$\bZ \times \bZ$ bigrading from the two loop directions.

\subsection{The CY parameter constraint}

In the additive parametrization $q_i = e^{h_i}$, the three
equivariant parameters of the $\Omega$-background satisfy the CY
condition:
\[
 h_1 + h_2 + h_3 = 0
 \qquad\text{equivalently}\qquad
 q_1 q_2 q_3 = 1.
\]
Two independent parameters survive.  The standard
$(q,t)$-parametrization identifies $q_1 = q$, $q_2 = t^{-1}$,
$q_3 = t/q$.

\begin{definition}[CY reflection identity]
\label{def:v3-qg-cy-reflection}
The \emph{trigonometric structure function} of the quantum toroidal
algebra is
\[
 G(x;\, q_1, q_2, q_3)
 = \frac{(1 - q_1 x)(1 - q_2 x)(1 - q_3 x)}
 {(1 - q_1^{-1}x)(1 - q_2^{-1}x)(1 - q_3^{-1}x)}.
\]
The \emph{CY reflection identity} is
\[
 G(x) \cdot G(1/x) = 1.
\]
\end{definition}

\begin{proposition}[CY reflection characterizes the CY condition]
\label{prop:v3-qg-cy-reflection}
Writing $q_1 q_2 q_3=\omega$, one has
\[
G(x)G(x^{-1})=\omega^2.
\]
Thus the connected exponential Calabi--Yau branch
$q_1q_2q_3=1$ implies the reflection identity
$G(x)G(x^{-1})=1$. Algebraically the branch $\omega=-1$ also
satisfies the squared identity; the converse requires the chosen
connected/formal branch.
\end{proposition}

\begin{proof}
Direct computation.  The numerator of $G(x) \cdot G(1/x)$ is
$\prod_i (1 - q_i x)(1 - q_i/x) = \prod_i (1 - q_i x)(x - q_i)/x$.
The denominator is
$\prod_i (1 - q_i^{-1}x)(1 - q_i^{-1}/x) = \prod_i
(1 - q_i^{-1}x)(x - q_i^{-1})/x$.
The ratio simplifies to $\prod_i (q_i/q_i^{-1}) = \prod_i q_i^2
= (q_1 q_2 q_3)^2$.
\end{proof}

\subsection{The Miki automorphism}

The Miki $S_3$ symmetry is an automorphism of the completed
quantum toroidal/DIM algebra. On parameters it permutes
$(q_1,q_2,q_3)$, but the automorphism is more than a cyclic
relabelling of the positive half. It uses the completed algebra,
Cartan sector, and opposite directions.

\begin{remark}[The Miki automorphism and the operadic action]
\label{rem:v3-qg-miki-non-operadic}
The $\Ethree$-operad acts on the derived center $\Zder(A_\cC)$
through the higher Deligne conjecture.  The Miki automorphism acts
on the quantum toroidal algebra itself, permuting the three
equivariant directions.  These are different actions at different
categorical levels: the $\Ethree$-action is on the bulk algebra
(the derived center), while the Miki automorphism is on the
boundary algebra (the chiral algebra itself).  The relationship
between them is mediated by the $\En$ operadic circle, but the
Miki automorphism requires the completed toroidal algebra.
\end{remark}

\subsection{The DDCA--toroidal bridge}

The deformed double current algebra
$\mathrm{DDCA}_k(\fgl_K) = U(\fgl_K \otimes \bC[u, v])$ carries
two spectral parameters from the two holomorphic directions of
$\bC^2$ in the $5$d theory.  The quantum toroidal algebra
$U_{q,t}(\widehat{\widehat{\fgl}}_K)$ carries two deformation
parameters from the two loop directions.

\begin{conjecture}[DDCA--toroidal bridge]
\label{conj:v3-qg-ddca-toroidal}
Writing $q = e^{\hbar_1}$, $t = e^{\hbar_2}$, the DDCA is the
rational ($\hbar_1, \hbar_2 \to 0$) degeneration of the quantum
toroidal algebra in the category of filtered algebras.  The two
one-parameter degenerations produce a commuting square:
\begin{center}
\begin{tikzcd}[column sep=1.8cm, row sep=1cm]
 U_{q,t}(\widehat{\widehat{\fgl}}_K)
 \arrow[r, "t \to 1"] \arrow[d, "q \to 1"']
 & Y(\widehat{\fgl}_K) \otimes U(\bC[v])
 \arrow[d, "\hbar_1 \to 0"] \\
 U(\bC[u]) \otimes Y(\widehat{\fgl}_K)
 \arrow[r, "\hbar_2 \to 0"']
 & \mathrm{DDCA}_k(\fgl_K)
\end{tikzcd}
\end{center}
The diagonal $\hbar_1 = \hbar_2 \to 0$ recovers the affine Yangian
$Y(\widehat{\fgl}_K)$ as the common degeneration.  The modular
characteristic is continuous across the bridge:
$\kappa_{\mathrm{ch}}(\mathrm{DDCA}_k(\fgl_K))
= K(k+K)/2 = \lim_{\hbar_i \to 0}
\kappa_{\mathrm{ch}}(U_{q,t})$.
\end{conjecture}


%%% ================================================================
%%% 6. DRINFELD CENTER
%%% ================================================================

\section{The Drinfeld center closes the $\En$ circle}
\label{sec:drinfeld-center}

The Drinfeld center $\cZ(\Rep^{\Eone}(A))$ is the universal
braided monoidal category receiving a central monoidal functor from
$\Rep^{\Eone}(A)$.  For the CoHA of a toric CY$_3$-fold:

\begin{theorem}[Drinfeld center for constructed CY$_3$ chiral loci]
\label{thm:v3-qg-drinfeld-center}
For $A_\cC$ in the constructed BZFN/CoHA loci of a toric
CY$_3$-fold~$X$:
\[
 \cZ(\Rep^{\Eone}(A_\cC))
 \;\simeq\;
 \Rep^{\Etwo}(\Zder(A_\cC))
\]
where $\Zder(A_\cC) = C^\bullet_{\mathrm{ch}}(A_\cC, A_\cC)$ is
the chiral derived center.  The forgetful functor
$U \colon \cZ(\Rep^{\Eone}(A_\cC)) \to \Rep^{\Eone}(A_\cC)$ is
the categorified averaging map.
\end{theorem}

\begin{remark}[Ordinary center and derived center]
\label{rem:v3-qg-ordinary-vs-derived}
For $A = H_k$ (Heisenberg at level~$k$): the ordinary center
$Z(A) = \bC$ has dimension~$1$; the derived center
$\Zder(A)$ has dimension~$3$, capturing $\Ext^1$ and $\Ext^2$
invisible to the commutant.  The five invariants matching at six
levels (conjectured for Heisenberg; proved end-to-end for $\bC^3$
via the CoHA) confirm that the Drinfeld center of the
$\Eone$-monoidal representation category is the correct
categorification of the chiral derived center.
\end{remark}

The passage $\Eone \to \Etwo$ via the Drinfeld center is the CY
incarnation of the $\Eone$-to-$\Etwo$ lift of Volume~I.  The
averaging map $\mathrm{av} \colon \fg^{\Eone} \to
\fg^{\mathrm{mod}}$ projects the ordered $R$-matrix to the scalar
$\kappa_{\mathrm{ch}}$; the Drinfeld center is the categorical
analogue, recovering the $R$-matrix data as the braiding on the
center category.

Direct restriction from $E_d$ to $\Etwo$ (for $d \geq 3$) gives
only symmetric braiding:
\[
\pi_1(\Conf_2(\bR^d))=0,\qquad
\pi_1(\Conf_2(\bR^d)/\mathfrak S_2)=\bZ/2
\quad(d\ge3).
\]
The Drinfeld center is the source of the nontrivial centre-side
half-braiding for CY$_3$ quantum groups.


%%% ================================================================
%%% 7. SHADOW = GW(C^3)
%%% ================================================================

\section{Shadow = $\GW(\bC^3)$ and wall-crossing}
\label{sec:shadow-gw}

\subsection{The shadow tower as Gromov--Witten theory}

At $\kappa_{\mathrm{ch}} = \Psi$ (the deformation parameter of the
low-spin defect channel), the shadow obstruction tower recovers the
genus-$1$ constant-map contribution for $\bC^3$. The all-genus virtual
identification requires the uniform-weight and $\Phi_3$ hypotheses;
the motivic identification remains conjectural.

\begin{proposition}[Genus-one shadow and the DT character of $\bC^3$]
\label{prop:v3-qg-shadow-gw}
At genus~$1$, the shadow specialization $\kappa_{\mathrm{ch}}=\Psi$
matches the constant-map/DT genus-one channel for $\bC^3$. On the DT
side, the MacMahon function
$M(q) = \prod_{n \geq 1}(1 - q^n)^{-n}$ encodes the partition
function as the second-quantized bar/DT character in the non-compact
degree-zero normalization.
\end{proposition}

The identification shadow = GW operates at three levels:
\begin{enumerate}[(i)]
 \item \emph{Genus-$1$}: the shadow Eisenstein theorem
 $D_2(\cA, s) = -24\kappa_{\mathrm{ch}} \cdot \zeta(s)\zeta(s-1)$
 (Volume~I) gives the genus-$1$ DT partition function via the
 Borcherds lift.
 \item \emph{Virtual}: under the uniform-weight and $\Phi_3$
 hypotheses, the shadow coefficient $S_r$ at degree~$r$ matches the
 virtual Euler characteristic of the moduli of $r$-dimensional sheaves.
 \item \emph{Motivic}: the full chiral quantum group structure
 ($R$-matrix, coproduct, RTT) is conjecturally matched with the
 motivic DT invariants of Kontsevich--Soibelman.
\end{enumerate}

\subsection{Wall-crossing as MC gauge equivalence}

For a toric CY$_3$ with fan~$\Sigma$, Bridgeland stability
conditions partition the space of central charges into chambers.
Each chamber $\sigma$ determines a quiver presentation
$(Q_\sigma, W_\sigma)$ and a CoHA
$\CoHA(Q_\sigma, W_\sigma)$.

\begin{conjecture}[Wall-crossing = MC gauge equivalence]
\label{conj:v3-qg-wall-crossing-mc}
Wall-crossing between stability chambers $\sigma$ and $\sigma'$ is
Maurer--Cartan gauge equivalence in the modular convolution algebra:
$\Theta_\sigma \sim \Theta_{\sigma'}$.  The Kontsevich--Soibelman
wall-crossing formula is the MC gauge transformation.  The global
chiral algebra is the homotopy colimit over the chamber
decomposition:
$A_\cC = \operatorname{hocolim}_\sigma A_{Q_\sigma}$.
\end{conjecture}

For $\bC^3$ (single chamber): the conjecture is vacuous.  For the
resolved conifold (two chambers): the wall-crossing formula is the
flop relation on DT partition functions (Bryan--Steinberg), and the
MC gauge equivalence is the bar-cobar transport across the flop.

\begin{remark}[The four faces of wall-crossing]
\label{rem:v3-qg-four-faces-wc}
Wall-crossing admits four descriptions in the programme:
\begin{enumerate}[(i)]
 \item \emph{DT/BPS}: change of BPS spectrum across a wall of
 marginal stability.
 \item \emph{MC gauge}: the universal Maurer--Cartan element
 $\Theta_A$ transforms by gauge equivalence in the modular
 convolution algebra.
 \item \emph{Scattering diagram}: the Gross--Siebert consistency
 condition IS the $\Eone$ Maurer--Cartan equation (proved for
 affine KM).
 \item \emph{Root system}: reflections in the Weyl group of the
 BKM superalgebra $\fg_X$.
\end{enumerate}
The identification of (ii) with (iii) is proved; the identification
with (i) and (iv) is conjectural in the CY$_3$ setting.
\end{remark}


%%% ================================================================
%%% 8. CY_4 and the p_1-twisted double current algebra
%%% ================================================================

\section{The $\CY_4$ obstruction and the $p_1$-twisted double
current algebra}
\label{sec:cy4-p1-twisted}

The fourth-dimensional stage of the holomorphic Chern--Simons
hierarchy is obstructed.  The untwisted quantum group associated to a
CY$_4$ geometry is replaced by a $p_1$-twisted double current algebra:
the first Pontryagin class of the tangent bundle appears as a cocycle
on the Drinfeld double.

\subsection{The obstruction in $\pi_4(BU)$}

\begin{theorem}[CY$_4$ first obstruction]
\label{thm:cy4-obstruction-pi4-BU}
Let $X$ be a compact CY$_4$ manifold. The obstruction to promoting
$\Phi(D^b\Coh(X))$ to an $E_4$-algebra lies in
$\pi_4(BU) \cong \bZ$, realised as the framing class of the stable
normal bundle to the cyclic trace on $X$.
\end{theorem}

\begin{proof}
By Bott periodicity,
$\pi_4(BU) = \pi_3(U) = \bZ$.
The equality $\pi_3(BU)=0$ follows from $\pi_2(U)=0$.  The
$E_1$-direction framing of the cyclic trace on $X$ carries a
characteristic obstruction class in $H^4(X;\bZ)$, represented by
$p_1(T_X)$.  Scalar obstructions are obtained by pairing this class
with specified four-cycles, or by the degree-compatible invariant
$\int_X p_1(T_X)^2$ on the real eight-manifold~$X$.
Obstruction theory for $E_1$-to-$E_4$ promotion locates this first
class in the $\pi_4(BU)$ stage. Its vanishing removes only this
obstruction; a genuine $E_4$ lift still requires the higher
obstruction tower and holomorphic-twist, framing, and completion
witnesses.
\end{proof}

\begin{remark}
Contrast $\CY_3$: there $\pi_3(BU) = 0$, so this first scalar
obstruction vanishes. The CY$_4$ case is the first step where this
homotopy group is non-trivial; the full higher obstruction tower
remains separate.
\end{remark}

\subsection{The Drinfeld double presentation}

\begin{definition}[$p_1$-twisted double current algebra]
\label{def:p1-twisted-double-current}
Let $X$ be a CY$_4$ and $\bg$ a simple Lie algebra.
Let $Y^+(\bg)$ and $Y^-(\bg)$ be the positive and negative
half-Yangians (in the Drinfeld realisation: Chevalley generators
$\{x_i^{\pm,r}\}_{r \ge 0}$ subject to the Yangian relations over
$\bC[\hbar]$).
The \emph{$p_1$-twisted double current algebra} associated to $X$ is
the Drinfeld double
\[
Y_{p_1}(X) \defeq Y^+(\bg) \bowtie_{c} Y^-(\bg)
\]
equipped with the cocycle
\[
c(x, y) \;=\; \frac{1}{24}\,
\Big\langle\, x \cup y \cup p_1(T_X),\ [X] \,\Big\rangle,
\qquad x,y \in H^{\bullet}(X;\bg),
\]
where $p_1(T_X) \in H^4(X;\bZ)$ is the first Pontryagin class and
$[X] \in H_8(X;\bZ)$ the fundamental class.
\end{definition}

\begin{remark}
The algebra $Y_{p_1}(X)$ is a $p_1$-twisted Drinfeld double.  Its
cocycle depends on $X$ through $p_1$, and its RTT presentation contains
a genuine $O(\hbar^2)$ correction absent from ordinary Yangians; the
notation ``CY$_4$ Yangian'' suppresses the obstruction.
\end{remark}

\subsection{The RTT relation through $O(\hbar^2)$}

\begin{conjecture}[$p_1$-twisted RTT]
\label{conj:cy4-rtt-p1}
Let
\[
N_{p_1^2}(X)\defeq \int_X p_1(T_X)^2 .
\]
The $R$-matrix of $Y_{p_1}(X)$ satisfies
\[
R_{p_1}(u) \;=\; 1 \;+\; \frac{\hbar\,P}{u}
\;+\; \hbar^2\,\frac{N_{p_1^2}(X)\,\Omega_{\mathrm{Cartan}}}{u^2}
\;+\; O(\hbar^3),
\]
where $P$ is the permutation operator and $\Omega_{\mathrm{Cartan}}$
the Cartan component of the quadratic Casimir.
\end{conjecture}

\begin{remark}[Proof obligation]
\label{rem:cy4-rtt-p1-proof-obligation}
The displayed $O(\hbar^2)$ term requires a deformation-complex
calculation for the $p_1$-twisted double, or a primary source proving
the same RTT normalisation.  Whitehead's lemma supplies only the
cohomological vanishing input.
\end{remark}

\subsection{CY$_4$ twisted-supergravity analogue}

\begin{conjecture}[CY$_4$ twisted-supergravity boundary algebra]
\label{conj:cy4-twisted-supergravity-boundary-algebra}
The boundary algebra of the conjectural CY$_4$ twisted-supergravity
analogue on $X \times \bR$, with $X$ a CY$_4$, is the $p_1$-twisted
double current algebra $Y_{p_1}(X)$.
\end{conjecture}

\begin{remark}[Evidence]
The statement is the CY$_4$ analogue of the 5d hCS/Yangian-VOA
theorem on $\bR\times\bC^2$ and of the 6d hCS framework of
Section~\ref{sec:hcs-hierarchy}. Since a CY$_4$ has real dimension
8, the space $X\times\bR$ is 9-dimensional. The statement is a
CY$_4$ twisted-supergravity conjecture. The cocycle $c$ is
predicted by the one-loop anomaly, proportional to $p_1(T_X)$.
\end{remark}

\subsection{Comparison with Maulik--Okounkov}

\begin{proposition}[Maulik--Okounkov comparison]
\label{prop:cy4-MO-comparison}
At $X = \bC^4$, the twist $c$ degenerates to the equivariant
Omega-background of Maulik--Okounkov, and the specialisation map
\[
Y_{p_1}(\bC^4)\ \longrightarrow\ U_{q,t_1,t_2}(\widehat{\widehat{\gl}}_1)
\]
identifies the $p_1$-twisted double current algebra with the
quantum toroidal $\gl_1$ algebra of Maulik--Okounkov.
\end{proposition}

\begin{remark}[Attribution]
Cited from Maulik--Okounkov, \emph{Quantum groups and quantum
cohomology}, Ast\'erisque 408 (2019), and Schiffmann--Vasserot.
The vanishing of $N(\bC^4)$ on the non-compact target is cancelled by
the equivariant integration that produces the $(q,t_1,t_2)$
parameters; this is the standard refinement dictionary.
\end{remark}

\subsection{$p_1$-twisted stratum of $K3 \times K3$}

\begin{proposition}[$K3 \times K3$ is $p_1$-twisted]
\label{prop:k3k3-p1-twisted}
For $X = K3 \times K3$,
\[
N_{p_1^2}(X) \;=\; \int_{K3 \times K3} p_1(T_X)^2
\;=\; 4608,
\]
and $p_1(T_X)\neq0$. Thus $K3\times K3$ lies on the $p_1$-twisted
stratum.  The untwisted Drinfeld double and pointwise $E_4$ lift
require separate vanishing data.
\end{proposition}

\begin{proof}
The Pontryagin class splits: $p_1(T_{K3 \times K3}) = \pi_1^* p_1(T_{K3})
+ \pi_2^* p_1(T_{K3})$, where $\pi_i$ are the two projections.
For a $K3$ surface $p_1(T_{K3}) = -2 c_2(T_{K3}) = -48\,[\pt]$ in
$H^4(K3)$. Therefore
\[
p_1(T_{K3\times K3})
=-48(\pi_1^*[\pt]+\pi_2^*[\pt]).
\]
Squaring gives
\[
p_1(T_{K3\times K3})^2
=2\cdot48^2\,\pi_1^*[\pt]\,\pi_2^*[\pt],
\]
and integration over $K3\times K3$ gives $2\cdot48^2=4608$.
\end{proof}

\begin{remark}
The product $K3\times K3$ is off the $S^4$-framed stratum. Its
natural CY$_4$ quantum group is $p_1$-twisted.
\end{remark}


%%% ================================================================
%%% Codimension-2 defects, CoHA positivity, and the toroidal formal disk
%%% ================================================================

\section{Codimension-$2$ defects, $\CoHA$, and the toroidal formal disk}
\label{sec:6dhCS-defect-coha-toroidal}

The 6d holomorphic Chern--Simons hierarchy has three load-bearing
tests: the codimension-$2$ defect algebra on a curve
$C \subset \bC^3$, the identification
$\CoHA(\bC^3) \simeq Y^+(\widehat\fgl_1)$ as the \emph{positive half}
of $\cW_{1+\infty}$, and the formal-disk specialisation of the
toroidal algebra via the Schiffmann--Vasserot $\CoHA$ and the Miki
$DIM$ presentation.

\begin{theorem}[Codim-$2$ low-spin defect at non-critical level]
\label{thm:codim2-defect}
Let $Y = \bC^3$ with holomorphic volume form
$\Omega_Y = dz_1\wedge dz_2\wedge dz_3$ and let
$C=\{z_2=z_3=0\}\subset Y$ be a codim-$2$ holomorphic curve. The
6d holomorphic Chern--Simons defect on~$C$ has a Heisenberg current
with level
\[
\Psi_{\mathrm{lev}}=-\sigma_2(h_1,h_2,h_3),
\]
where $h_1$ is the tangent weight and $(h_2,h_3)$ are the normal
weights, subject to $h_1+h_2+h_3=0$. The Sugawara stress tensor has
central charge $c=1$.
\end{theorem}

\begin{proof}
Reduction of 6d hCS along $N_{C/Y}$ is a codim-$2$ boundary
condition; the descent of the gauge algebra from $\mathrm{hol}(Y)$
to $\mathrm{hol}(C)$ factors through the normal Koszul complex,
producing a chiral extension of $\widehat\fgl_1$. The equivariant
normal parameters are $h_2,h_3$ and the tangent weight is $h_1$.
Their elementary symmetric polynomial $\sigma_2$ gives the
Heisenberg level $\Psi_{\mathrm{lev}}=-\sigma_2$. The vanishing
$\sigma_1=0$ is the CY condition. Non-criticality
$\Psi_{\mathrm{lev}}\ne0,\infty$ places the Heisenberg current on
the Sugawara locus.
\end{proof}

\begin{theorem}[$\CoHA(\bC^3)$ as the positive half]
\label{thm:coha-positive-half}
The critical cohomological Hall algebra of $\bC^3$
($=$~Jordan quiver with cubic potential) is isomorphic to the
positive half $Y^+(\widehat\fgl_1)$ of the affine Yangian of
$\widehat\fgl_1$:
\begin{equation}\label{eq:coha-positive}
\CoHA(\bC^3) \;\simeq\; Y^+(\widehat\fgl_1).
\end{equation}
The full $\cW_{1+\infty} \simeq Y(\widehat\fgl_1)$ is recovered as
the Drinfeld double $Y^+ \bowtie Y^-$, with the opposite half
$Y^-$ realised either by the opposite orientation $(\bC^3)^{\op}$
(conjugate quiver with potential $-W$) or by a Koszul dual
$\CoHA^!$ construction.
\end{theorem}

\begin{proof}
Schiffmann--Vasserot prove $\CoHA(Q, W) \simeq Y^+$ as the Hopf
algebra generated by simple-object fundamental classes, with Hall
multiplication for positive roots; the shuffle presentation exhibits
only $Y^+$. The full Yangian with Cartan currents and negative
roots requires either a second copy via $(\bC^3)^{\op}$ or the
derived Drinfeld double. Thus the precise identification is
$\CoHA \simeq Y^+$. Character coincidence with
$\cW_{1+\infty}$ is weaker than a bar equivalence.
\end{proof}

\begin{theorem}[Toroidal formal-disk identification]
\label{thm:toroidal-formal-disk}
On the formal disk $D = \Spec\Bbbk[[z]]$, the completed quantum
toroidal algebra $U_{q,t}(\widehat{\widehat\fgl}_1)$ carries the
structure of an $\Eone$-chiral quantum group: (i)~the completed Miki
$S_3$-automorphism realizes the $\Omega$-background Weyl symmetry on
the branch $q_1 q_2 q_3 = 1$; (ii)~the spectral $R$-matrix from the
completed SV/toroidal structure satisfies the additive Yang--Baxter equation;
(iii)~the formal-disk chiral quantum-group equivalence holds at that
level with associator dependence controlled by the Etingof--Kazhdan
functor on the Koszul locus.
\end{theorem}

\begin{remark}[Global toroidal: conditional]
\label{rem:toroidal-global}
The extension of Theorem~\ref{thm:toroidal-formal-disk}
from the formal disk to a global curve $\bP^1 \times \bP^1$ is
conditional on chain-level topologization of the underlying
class~M chiral algebra: the formal-disk statement is unconditional,
the global statement inherits the class~M chain-level
topologization hypothesis and the modular-family base-change
problem (proved cohomologically;
the chain-level statement depends on Francis--Gaitsgory base change
for the relative Ran prestack).
\end{remark}

\begin{remark}[Geometric vs algebraic model discipline]
\label{rem:geom-vs-alg}
The parameters $(h_1,h_2,h_3)$ are formal algebraic variables in
$\End^{\mathrm{ch}}_\cA$. Fulton--MacPherson variables on
$\FM_k(\bC)$, the Hopf-algebra identity
$\CoHA(\bC^3)\simeq Y^+$, and BZFN centre comparisons live in three
distinct ambients: configuration-space geometry, algebraic Hopf
algebras, and the pair $(\cA^{\mathrm{ch}},\cA_{\mathrm{mode}})$.
\end{remark}

\begin{remark}[Proof ambients]
\label{rem:6d-proof-ambients}
Theorem~\ref{thm:codim2-defect} rests on the formal BRST calculation
at non-critical level (\emph{proved cohomologically for affine KM;
class~M cases via DS reduction conditional}). Theorem~\ref{thm:coha-positive-half}
uses the Schiffmann--Vasserot positive-half normalization. Identifying
the CoHA directly with $\cW_{1+\infty}$ is only a character-level
shortcut unless the Drinfeld double is included.
Theorem~\ref{thm:toroidal-formal-disk} is \emph{proved}
at the formal disk; global extension conditional per
Remark~\ref{rem:toroidal-global}.
\end{remark}


%%% ================================================================
%%% REFERENCES
%%% ================================================================

\section*{References}

\begin{enumerate}[label={[\arabic*]}]
\item
 R.~Lorgat,
 \emph{Modular Koszul Duality},
 Volume~I of the Modular Koszul Duality Programme, 2026.

\item
 R.~Lorgat,
 \emph{$\Ainf$ Chiral Algebras and 3D Holomorphic-Topological QFT},
 Volume~II of the Modular Koszul Duality Programme, 2026.

\item
 R.~Lorgat,
 \emph{Calabi--Yau Categories, Quantum Groups, and BPS Algebras},
 Volume~III of the Modular Koszul Duality Programme, 2026.

\item
 R.~Lorgat,
 \emph{The Calabi--Yau-to-chiral functor and the
 $\kappa_\bullet$-spectrum},
 standalone paper (companion to this paper), 2026.

\item
 K.~Costello,
 \emph{Holomorphic Chern--Simons theory and the super Yangian},
 preprint, 2017.

\item
 K.~Costello, E.~Witten, and M.~Yamazaki,
 \emph{Gauge theory and integrability, I, II},
 Notices ICCM \textbf{6} (2018), 46--119 and 120--149.

\item
 O.~Schiffmann and E.~Vasserot,
 \emph{Cherednik algebras, $\cW$-algebras and the equivariant
 cohomology of the moduli space of instantons on $\mathbb{A}^2$},
 Publ.\ Math.\ IHES \textbf{118} (2013), 213--342.

\item
 M.~Rapcak, Y.~Soibelman, Y.~Yang, and G.~Zhao,
 \emph{Cohomological Hall algebras, vertex algebras, and
 instantons},
 Comm.\ Math.\ Phys.\ \textbf{376} (2020), 1803--1873.

\item
 T.~Prochazka and M.~Rapcak,
 \emph{$\cW$-algebra modules, free fields, and Gukov--Witten
 defects},
 JHEP \textbf{2019}(5), 159.

\item
 K.~Costello and O.~Gwilliam,
 \emph{Factorization Algebras in Quantum Field Theory},
 Cambridge University Press, 2017, 2021.

\item
 D.~Maulik, N.~Nekrasov, A.~Okounkov, R.~Pandharipande,
 \emph{Gromov--Witten theory and Donaldson--Thomas theory, I, II},
 Compos.\ Math.\ \textbf{142} (2006), 1263--1285 and 1286--1304.

\item
 M.~Kontsevich and Y.~Soibelman,
 \emph{Stability structures, motivic Donaldson--Thomas invariants
 and cluster transformations},
 preprint arXiv:0811.2435, 2008.
\end{enumerate}

\end{document}
